% Chapter-specific authoring template generated by Mission 13.
% Target chapter file: manuscript/chapters/08-social-network.tex
% Do not paste substantive findings without checking mapped source files.

\chapter{Witnesses, Contacts, and Social Network}
\label{chap:08-social-network}

\classification{Evidence status: <set during drafting>. Confidence baseline: Moderate.}

\section{Purpose}
Present witness evidence, contact leads, relationship map, and reliability constraints without assigning motive.

\section{Scope}
Use only the mapped source files and related registers. Do not add new research unless a critical evidentiary defect is discovered and logged.

\section{Evidence Summary}
Inputs: witness index; cross-reference map; relationship map; person register; transcription queue.

Source files: \texttt{research/indexes/witness_index.csv; research/indexes/cross_reference_map.csv; research/victimology/relationship_map.csv; research/database/person_register.csv}

\section{Key Findings To Address}
Witnesses support discovery and some evening sightings, but identity continuity and memory limits are material; Thomson/Boxall/Webb relationship claims require careful source separation.

\section{Body}
<Draft evidence-led chapter prose here. Separate established fact, supported inference, hypothesis, context and unknowns.>

\section{Tables}
Planned tables: witness_index_ready_for_latex.tex; witness reliability table; relationship map table.

\section{Figures}
Planned figures: Network diagram optional; no family/person photographs until rights-cleared and evidentiary relevance is established.

\section{Outstanding Questions}
Manual transcription of priority pages; original police statements; Thomson interview notes; full Boxall/Jestyn source trail

\section{Confidence Statement}
Baseline confidence: Moderate. Explain changes only if drafting exposes a documented source limitation.

\section{References}
Use end references from \texttt{bibliography/master_references.bib} and source/register IDs.
